\section{Introduction}\label{sec:intro}

A cell restores a membrane gradient after a perturbation. A learning agent revises a state
abstraction after a failed prediction. A court amends a procedure after a class of disputes has
defeated the old one. In each case one may ask a structural question that does not depend on the
material: can the repair be credited to the system, or was it supplied from outside? It can be
credited to the system if the defect is recorded in the system's own state, the repair is chosen from
the system's own repertoire and charged to its own budget, the repair step is one of the system's own
transitions, and the result is checked by the system's own instrument. If any of these records or
steps belongs to an experimenter, an engineer or an external administrator, the repair has been
supplied.

The earlier papers of this series give exact tools for supplied repair. \FoundIV{} collects
fifty-two structural laws about fixed data: when a readout descends through a quotient, how a
failed descent is repaired by refining the source or coarsening the target, when a residual must be
priced \citep{Tsiokos2026FoundIV}. In all of them the repair is supplied by the theorist: a map
fails to descend and the theorist refines its source. This paper asks what the same operations look
like when they are coordinates of the system's own state. The basic move is
\[
  O \ \longmapsto\ \Carried_{\mathbb S}(O),
\]
where a record is \emph{carried} by a system $\mathbb S$ if it is a coordinate of the state of
$\mathbb S$, is updated only by transitions of $\mathbb S$, and was produced by a trajectory of
$\mathbb S$ rather than by a fallback or an external edit (\Cref{def:d3}). Carriedness is a
condition on the source of a record. It attributes no purpose to the system.

\subsection{A worked example: a repair the system performs}\label{sec:intro-example}

The smallest system that repairs itself has two states (\Cref{fig:device}). State $0$ is damaged and state $1$ is
repaired; the transition kernel allows a move from either state to $1$, and the system's trajectory
is $0, 1, 1, \dots$. The system carries three things in its state. An \emph{instrument} reports a
defect exactly when the state is $0$, permits a repair of the first structural sort (P1, a
lower-operator rewrite) exactly at $0$, and accepts a state on re-audit exactly when it is $1$. A
\emph{ledger} has one budget line. A \emph{repair generator} sends each defect record to a P1 move
charged to that line. Every record has the source tag ``committed state'' and is produced by the
system's own trajectory.

The first step of the trajectory, $0 \to 1$, is a lawful repair step in the sense of
\Cref{def:d4-step}: the defect is detected, the gate permits the move, the move is charged to the
carried ledger and is admissible at $0$, the kernel allows $0 \to 1$, and the re-audit at $1$ passes.
The system's history records one challenge episode at time $0$ with three entries --- the defect,
the move and the audit --- read at time $1$. The record that the system's apparatus was reinstated
at time $0$ is linked to that episode, uses its defect and audit, and names $\tau(0)$ and $\tau(1)$
as its source and target (\Cref{thm:e3-episode}). A history with no episode supports no such link.

Give the device two instrument levels. The lower level carries its own instrument record. The upper
level checks the lower one --- is the lower level's record among the records its instrument makes
visible? --- and stores the answer, which is ``yes'', as a carried audit record. The records counted
at the two levels, $\{i_0\}$ and $\{i_1, a\}$, are disjoint and lie in a record universe of size
$3$, so the tower's footprint is at most its capacity, $1 + 2 \le 3$ (\Cref{thm:e2-tower}).

\begin{figure}[t]
\centering
\begin{tikzpicture}[>=Stealth,
  st/.style={circle, draw, minimum size=9mm},
  rec/.style={rectangle, draw, rounded corners=2pt, font=\small, inner sep=3pt},
  lab/.style={font=\footnotesize}]
  \node[st] (z0) at (0,0) {$0$};
  \node[st] (z1) at (3.6,0) {$1$};
  \draw[->, thick] (z0) -- node[above=2pt, lab] {repair (P1 move)} node[below=2pt, lab] {$\tau(0)\to\tau(1)$} (z1);
  \draw[->] (z1) edge[loop right] node[lab, right] {stay} (z1);
  \node[lab, below=0.35cm of z0, align=center] {defect detected\\gate open};
  \node[lab, below=0.35cm of z1, align=center] {re-audit passes};
  \node[rec] (l1) at (8.6,0.7) {level $1$: $\{i_1, a\}$};
  \node[rec] (l0) at (8.6,-0.7) {level $0$: $\{i_0\}$};
  \draw[->] (l1) -- (l0);
  \node[draw, dashed, rounded corners, fit=(l0)(l1), inner sep=6pt] (box) {};
  \node[lab, above=1pt of box] {record universe of size $3$};
  \node[lab, below=2pt of box, align=center] {level $1$ audits level $0$:\\is $i_0$ visible? $a = $ yes};
\end{tikzpicture}
\caption{The repair device of \Cref{def:device}. Left: its trajectory $0, 1, 1, \dots$ takes the
lawful repair step $0 \to 1$ (\Cref{thm:e3-episode}). Right: the two-level audit tower, whose counted
records fill a universe of size $3$ (\Cref{thm:e2-tower}).}\label{fig:device}
\end{figure}

The example is small on purpose: the defect, move, audit and ledger records of the device each have
one value, the tower adds two instrument and two audit values, and formedness and completeness of
the inventories are declared. We call such instances \emph{weak} (\Cref{sec:setting-status}).
What it shows is that the conditions of a self-credited repair --- carried defect, generated move,
own budget, own transition, own audit, linked history --- can be met together by an explicit system
and checked.

\subsection{How the catalog is sorted}\label{sec:intro-sorting}

The paper states sixteen entries, E1--E16, over six definitions, D1--D6, and sorts them by what is
proved (\Cref{tab:catalog}). A \emph{law} is a general statement together with an instance in which
its hard hypothesis is proved about a concrete system. Five entries qualify: E1 (perpetual
self-extension), E2 (bounded reflexivity), E3 (self-maintaining reclosure), E12 (individuation) and
E13 (repair transport). A \emph{general theorem} has ordinary hypotheses and nothing open; E6 (priced
exposure allocation) and E9 (probe acquisition) are general theorems, and E6 has a numerical
instance. The remaining entries are \emph{certificate specifications}: the Lean
statement shows that a complete certificate yields a classification, and supplying the certificate
for a real system is the open obligation. Two entries, E5 and E8, are best read as definitions with
consistency lemmas.

\subsection{Main results}\label{sec:intro-results}

Each result below is proved in the text and checked in Lean.

\begin{headline}[E1: perpetual self-extension]
A repair family that recurs beyond every bound, and installs a strict refinement whenever the
challenge is new, makes infinitely many strict self-extensions under a challenge that is always new
(\Cref{prop:e1-forcing}). A clocked device performs a lawful repair at every odd time and installs
pairwise distinct refinements that meet both conditions (\Cref{thm:e1-clocked}).
\end{headline}

\begin{headline}[E2: audit capacity]
For a declared finite inventory of instrument levels whose counted record sets are pairwise disjoint
subsets of a record universe of size $\capc(z)$, the total footprint is at most $\capc(z)$
(\Cref{thm:e2-declared}). The two-level audited tower of the device attains $3 \le 3$
(\Cref{thm:e2-tower}).
\end{headline}

\begin{headline}[E3: an executed, history-linked reinstatement]
Along the device's trajectory the step $0 \to 1$ is a lawful repair, and the reinstatement of its
apparatus is linked to a history episode dated at that step whose defect, move and audit entries are
carried (\Cref{thm:e3-episode}).
\end{headline}

\begin{headline}[E6: priced allocation]
Given budget and cap multipliers that satisfy stationarity and complementary slackness for a capped
linear objective, the allocation is a global optimum, including when it sits at a cap
(\Cref{thm:e6-caps}); with weights $(2, 1)$, caps $(1, 5)$ and budget $3$ the optimum $(1, 2)$ has a
positive cap multiplier (\Cref{thm:e6-atcap}). For any supplied multipliers, a selected probe with
positive marginal discharge forces a positive price (\Cref{thm:e6-kkt}).
\end{headline}

\begin{headline}[E12 and E13: individuals and transport]
Parts that are closed under repair form an intersection-closed family, and integrated individuals
are equal or disjoint (\Cref{prop:e12-structure}). In a three-state system the part made of a
damaged state and its repair is an integrated individual, while the healthy state alone passes every
test except repair closure (\Cref{thm:e12-repair}). Repair transports compose with added payments,
and a chain of $k$ transported repairs lowers the receiver's obstruction by at least $k$
(\Cref{thm:e13-chain}).
\end{headline}

\begin{headline}[E9 and E15: acquisition and reallocation]
Over a finite nonempty list of admissible acquisition candidates that is complete for admissibility,
some listed candidate maximizes discharge per cost among all admissible candidates
(\Cref{thm:e9}). If an offline phase gates external exchange to zero and the shared budget is
exhausted in both phases, the offline internal capacity equals the online internal capacity plus
the freed exchange share (\Cref{prop:e15-identity}).
\end{headline}

\begin{headline}[D1 and D2: repair joins and residual traces]
The repair join $q \vee r$ is the greatest lower bound of $q$ and $r$ in the refinement order and is
idempotent, commutative and associative up to fibre equivalence (\Cref{prop:d1}). If the adequacy
residual satisfies the chain rule under a new probe and the conditional term is positive
semidefinite, the residual trace does not increase (\Cref{prop:d2}).
\end{headline}

\subsection{Map of results}\label{sec:intro-map}

\Cref{tab:map} lists the numbered results. Each is proved after its statement, except \Cref{prop:e2-self}, whose source is given in \Cref{tab:provenance}.

\begingroup\small
\begin{longtable}{@{}>{\raggedright\arraybackslash}p{3.0cm}>{\raggedright\arraybackslash}p{12.4cm}@{}}
\caption{Map of results.}\label{tab:map}\\
\toprule
Result & What it establishes \\
\midrule
\endfirsthead
\toprule
Result & What it establishes \\
\midrule
\endhead
\bottomrule
\endfoot
\cref{prop:d1} & The repair join is, up to equality of fibres, the greatest lower bound in the refinement order, and it is idempotent, commutative and associative. \\
\cref{prop:d2} & The conditional surplus is nonnegative exactly when the probe score is at most the quotient score; a positive semidefinite residual has nonnegative trace; under the chain rule with a positive semidefinite subtracted term, adjoining a probe does not increase the residual trace, and a probe $M=BL$ with $K_{DM\mid L}=0$ leaves it unchanged. \\
\cref{prop:device-step} & In the repair device, $0\to1$ is a lawful repair step. \\
\cref{thm:e3-episode} & In the repair device with its declared trajectory, $0\to1$ is a lawful repair step, the reinstatement record is history-linked with all three entries attached and carried, and the episode is dated at the reinstatement; in a history with no episodes no reinstatement record is history-linked. \\
\cref{prop:e2-self} & A carried self-audit claim whose subject and accepting instrument are the same level, with no bridge to a different level, is not accepted. \\
\cref{thm:e2-declared} & For every declared capacity measure the footprints of the declared levels sum to at most the capacity, and a level with the own record of a declared level at a different index is not declared. \\
\cref{thm:e2-tower} & The two-level audited tower on the repair device checks its lower level, forms a declared capacity measure with footprint $3\le3$, and keeps the lawful repair step $0\to1$. \\
\cref{thm:e6-kkt} & For a multiplier witness with positive marginal costs, a selected probe of positive discharge forces a positive price equal to the ratio of every selected probe, unselected probes have no larger ratio, and a slack budget forces price zero. \\
\cref{prop:e6-sufficient} & A multiplier witness at a feasible allocation, for a concave differentiable objective whose partial derivatives are the supplied discharges and with the recorded spend, gives an optimal allocation. \\
\cref{thm:e6-caps} & For a capped linear objective on two coordinates, a feasible allocation below its caps with multipliers satisfying the stated conditions is optimal. \\
\cref{thm:e6-atcap} & With weights $(2,1)$, caps $(1,5)$ and budget $3$, the allocation $(1,2)$ is optimal, with one coordinate at its cap; with unit weights it is also optimal. \\
\cref{prop:e1-forcing} & A repair family with an actual-family forcing certificate, facing a challenge that is new at every time, makes infinitely many strict self-extensions. \\
\cref{thm:e1-clocked} & The clocked device has a lawful repair step at every odd time and entries beyond every bound; it has an actual-family forcing certificate with a challenge new at every time, so it makes infinitely many strict self-extensions, with distinct refinements at distinct times. \\
\cref{prop:e12-structure} & In every individuation setting, intersections of repair-closed parts are repair-closed, and two integrated individuals are equal or disjoint. \\
\cref{thm:e12-repair} & On three states, $\{h,d\}$ is an integrated individual containing the damaged state and its repair, while $\{h\}$ is closed under steps and self-maintaining but not a candidate. \\
\cref{prop:e12-overlap} & Two overlapping maximal boundary regions on three states leave no integrated individual. \\
\cref{thm:e13-chain} & Repair transports compose, with payments adding, and along a repair transport a chain of $k$ source repairs lowers the receiver's obstruction by at least $k$. \\
\cref{thm:e13-cross} & An explicit repair transport exchanging two contexts lowers the receiver's obstruction along a two-step chain, and composing two transports gives one whose payment is the sum. \\
\cref{prop:e13-flat} & A paid, role-preserving map that carries every source repair to a receiver repair need not reduce any obstruction. \\
\cref{thm:e9} & Over a complete finite nonempty list of admissible acquisitions some listed acquisition maximizes discharge per cost among all admissible ones, also for lists mixing allocation increments and acquisitions, and a policy selecting a listed maximizer dominates every admissible alternative. \\
\cref{prop:e9-saturation} & Under the chain rule, a probe $M=BL$ with $K_{DM\mid L}=0$ discharges nothing. \\
\cref{prop:e15-identity} & When a shared budget binds and is exhausted in both phases, the offline internal share equals the online internal share plus the external share. \\
\cref{prop:e4-brittle} & A split pair for the update of a new challenge class through the minimal quotient of a compiled repair blocks descent of its readout. \\
\end{longtable}
\endgroup

\subsection{What is imported, and what is not claimed}\label{sec:intro-claims}

The Karush--Kuhn--Tucker conditions are standard \citep{BoydVandenberghe2004}; the chain rule and
same-family saturation of the adequacy residual, the split-pair descent criterion, the same-level
self-audit obstruction and the source-of-truth tags are taken from earlier papers of the series and
are not re-proved here \citep{Tsiokos2026Adequacy,Tsiokos2026FoundIV,Tsiokos2026FoundIII}. The
paper does not define life, cognition or institutions, and no biological, psychological or social
system is shown here to satisfy any law. The instances are small finite systems chosen to make each
hard hypothesis provable; the domain readings of \Cref{sec:scope} say which data an application
would have to supply, and they are not results.

\paragraph{Roadmap.} \Cref{sec:related} places the paper in the literature, and \Cref{sec:setting}
fixes the vocabulary, the six definitions and the catalog. \Cref{sec:concrete} states the five laws
with their instances, together with E6 and its numerical instance. \Cref{sec:theorems} gives the general theorems and \Cref{sec:templates}
the certificate specifications and definitions. \Cref{sec:scope} lists inferences the results do not
support and the domain readings, and \Cref{sec:lean} records what is formalized.
The supplement holds the source
concordance, the field lists of every certificate, the full register of non-implications, the finite
computations, the synthetic life-facing data, the axiom record and the replay commands.
